% 06-kiem-tra-gradient.tex - Finite-difference gradient checking
%
% Scope: two-sided finite difference formula; why it works; measured
% errors from the verification run; why this check is load-bearing.

\section{Kiểm tra gradient bằng sai phân hữu hạn}
\label{sec:ch01-gradient-check}

BPTT dễ cài đặt sai. Một dấu trừ nhầm chỗ, một chuyển vị thiếu, một chỉ số
\mintinline{python}{t} với \mintinline{python}{t+1} lẫn lộn; và gradient vẫn
là một ma trận trông hợp lý. Nó chỉ sai, và bạn sẽ không biết cho đến khi mô
hình không hội tụ, mà một mô hình không hội tụ thì có hàng tá nguyên nhân
khác ngoài gradient sai.

Cách duy nhất để biết chắc một cài đặt lan truyền ngược là đúng: so sánh nó
với gradient được tính bằng \textbf{sai phân hữu hạn}. Đây là kỹ thuật lâu đời
nhất trong tối ưu hóa số, và nó không phụ thuộc vào code của bạn; nó chỉ phụ
thuộc vào định nghĩa của đạo hàm.

\subsection{Công thức}

Với một tham số vô hướng $\theta$ và một hàm mất mát $L(\theta)$, đạo hàm
được xấp xỉ bằng:

\begin{equation}
  \frac{dL}{d\theta} \approx
  \frac{L(\theta + \varepsilon) - L(\theta - \varepsilon)}{2\varepsilon},
  \label{eq:finite-diff}
\end{equation}

với $\varepsilon$ là một số rất nhỏ. Đây là \textbf{sai phân trung tâm} bậc
hai: sai số tỉ lệ với $\varepsilon^2$, tốt hơn sai phân tiến bậc một (tỉ lệ
với $\varepsilon$). Cái giá là phải chạy lan truyền xuôi \emph{hai lần} cho
mỗi tham số; một lần với $+\varepsilon$, một lần với $-\varepsilon$.

Với một ma trận trọng số $\mat{W}$, ta áp dụng \eqref{eq:finite-diff} cho
từng phần tử một. Nếu $\mat{W}$ có $m \times n$ phần tử, ta cần $2mn$ lần
chạy lan truyền xuôi. Với RNN đồ chơi $3 \times 2$ của chúng ta, đó là 12
lần chạy cho $\mat{W}_{xh}$, 18 cho $\mat{W}_{hh}$, 6 cho $\mat{W}_{hy}$,
3 cho $\vect{b}_h$, và 2 cho $\vect{b}_y$; tổng cộng 82 lần chạy, chưa đầy
một phần nghìn giây trên CPU.

\subsection{Kết quả}

Tôi chạy với $\varepsilon = 10^{-5}$. Với mỗi tham số, tôi tính sai khác tuyệt
đối lớn nhất giữa gradient BPTT và gradient sai phân hữu hạn, rồi chia cho
tổng biên độ của hai gradient để có sai số tương đối.

\begin{measured}{Sai phân hữu hạn, $\varepsilon = 10^{-5}$}
\small
\begin{tabular}{@{}lrrr@{}}
\toprule
Tham số & $\max|\text{BPTT} - \text{FD}|$ & Sai số tương đối & Kết luận \\
\midrule
$\mat{W}_{xh}$ & $2.82 \times 10^{-11}$ & $1.35 \times 10^{-11}$ & Khớp \\
$\mat{W}_{hh}$ & $2.93 \times 10^{-11}$ & $2.28 \times 10^{-11}$ & Khớp \\
$\mat{W}_{hy}$ & $4.40 \times 10^{-11}$ & $8.69 \times 10^{-12}$ & Khớp \\
$\vect{b}_h$   & $4.88 \times 10^{-11}$ & $1.15 \times 10^{-11}$ & Khớp \\
$\vect{b}_y$   & $2.79 \times 10^{-11}$ & $4.07 \times 10^{-12}$ & Khớp \\
\bottomrule
\end{tabular}
\normalsize
\end{measured}

Tất cả gradient khớp với sai số tương đối dưới $1.4 \times 10^{-11}$. Đây là
độ chính xác ở mức float64: với $\varepsilon = 10^{-5}$, sai số của sai phân
trung tâm vào khoảng $\varepsilon^2 \cdot L''' \approx 10^{-10}$, cộng thêm
sai số làm tròn của phép cộng/trừ float64. Con số $10^{-11}$ nằm gọn trong
khoảng đó.

Nếu bạn thấy sai số tương đối trên $10^{-6}$ khi kiểm tra gradient trong dự
án của mình, có bug. Không có ngoại lệ.

\subsection{Vì sao việc này quan trọng}

Tôi biết một nghiên cứu sinh bỏ ba tuần debug một mô hình không hội tụ.
Nguyên nhân: gradient của $\mat{W}_{hh}$ bị thiếu một phép chuyển vị.
Một dòng code. Cả ba tuần đó có thể được tiết kiệm bằng 82 lần chạy sai
phân hữu hạn, tốn chưa đến một giây.

Mỗi chương trong cuốn sách này, trước khi in bất kỳ con số nào, đều chạy
gradient check trên code của chương đó. Kết quả được ghi vào research note
và không được chỉnh sửa. Đây không phải là thói quen tốt; đây là thứ duy
nhất ngăn tôi in một ma trận gradient sai lên giấy.

Trong thực tế, bạn không cần viết gradient check bằng tay như tôi làm ở đây.
PyTorch có \mintinline{python}{torch.autograd.gradcheck}, TensorFlow có
\mintinline{python}{tf.test.compute_gradient}. Nhưng hiểu thứ bên trong
những hàm đó; biết chính xác sai phân trung tâm là gì, $\varepsilon$ nên
chọn bao nhiêu, sai số kỳ vọng là bao nhiêu; là thứ phân biệt một kỹ sư
dùng được công cụ với một kỹ sư tin công cụ một cách mù quáng.
